\documentclass[10pt,a4paper]{article}

\usepackage[margin=0.1cm]{geometry}
\usepackage{multicol}
\usepackage{amsmath, amssymb}
\usepackage{tcolorbox}
\usepackage{enumitem}
\usepackage{makecell}
\usepackage{tabularx}
\usepackage{tikz}
\usepackage{pgfplots}
\pgfplotsset{compat=1.18} % latest stable

% compact lists
\setlist{nosep}

% compact sections
\setlength{\parindent}{0pt}
\setlength{\parskip}{0pt}

\begin{document}

\begin{center}
    {\Large \textbf{Estadística}} \\[4pt]
    \small % small subtitle here
\end{center}

\newcommand{\sectiontitle}[1]{\par\hrule\vspace{10pt}{\Large\bfseries #1}\par}

\sectiontitle{Conjuntos}

\tcbset{colback=gray!10, colframe=black, left=2pt, right=2pt, top=2pt, bottom=2pt, boxsep=2pt}

\begin{multicols}{2}
    \begin{tcolorbox}[title=Teoría de Conjuntos, equal height group=CONJ1]
        \begin{tabular}{rl}
            $x \in A$: & $x$ es un elemento de $A$ \\
            $A \subseteq B$: & todos los elementos de A están en B \\
            $|A| = \#(A)$: & número de elementos de A (cardinal) \\
            $\mathcal{P}(A)$: & conjunto de todos los subconjuntos de $A$ \\
            $|\mathcal{P}(A)| = 2^{|A|}$: & cardinal del conjunto de potencia \\
        \end{tabular}
    \end{tcolorbox}

    \begin{tcolorbox}[title=Función Característica, equal height group=CONJ1]
        Asigna 1 si el elemento está en $A$, 0 en caso contrario. Útil para operaciones booleanas. \\
        $ \chi_A(\omega) = \begin{cases} 1 & \text{si } \omega \in A \\ 0 & \text{si } \omega \notin A \end{cases}$
    \end{tcolorbox}
\end{multicols}

\begin{multicols}{2}
    \begin{tcolorbox}[title=Operaciones Básicas]
        \begin{tabular}{rl}
            \textbf{Intersección}: & $A \cap B = \{ x \in \Omega \mid x \in A \land x \in B \}$ \\
            \textbf{Unión}: & $A \cup B = \{ x \in \Omega \mid x \in A \lor x \in B \}$ \\
            \textbf{Diferencia}: & $A - B = \{x \in \Omega \mid x \in A \land x \notin B \}$ \\
            \textbf{Complemento}: & $A^c = \Omega - A = \{x \in \Omega \mid x \notin A \}$
        \end{tabular}
    \end{tcolorbox}
    \begin{tcolorbox}[title=Leyes de De Morgan]
        \begin{tabular}{rl}
            ${\left( A \cup B \right)}^c = A^c \cap B^c$ \\
            ${\left( A \cap B \right)}^c = A^c \cup B^c$
        \end{tabular}
    \end{tcolorbox}
    \begin{tcolorbox}[title=Propiedades adicionales]
        Dos conjuntos son disjuntos si $A \cap B = \emptyset$ \\
        $\emptyset^c = \Omega, \Omega^c = \emptyset$
    \end{tcolorbox}
\end{multicols}

\sectiontitle{Combinatoria}

\begin{multicols}{2}
    \begin{tcolorbox}[title=Número Binomial (Combinaciones sin Repetición), equal height group=COMB1]
        Número de subconjuntos de tamaño \( k \) en un conjunto de \( n \) elementos (orden no importa, sin repetición).
        \[
            \binom{n}{k} = C_k^n = \frac{n!}{k! \cdot (n-k)!}
        \] \\
        Ejemplo: \( \binom{7}{5} = 21 \) (equipos de 5 con 7 posibles jugadores).
    \end{tcolorbox}
    \begin{tcolorbox}[title=Número Binomial (Combinaciones sin Repetición), equal height group=COMB1]
        Número de formas de elegir \( k \) elementos de \( n \) tipos permitiendo repeticiones (multiconjuntos).
        \[
            CR_k^n = \binom{n + k - 1}{k} = \frac{(n + k - 1)!}{k! \cdot (n-1)!}
        \]
        Ejemplo: Repartir 12 caramelos en 4 personas: \( \binom{4 + 12 - 1}{12} = 455 \).
    \end{tcolorbox}
\end{multicols}

\begin{multicols}{2}
    \begin{tcolorbox}[title=Variaciones sin Repetición, equal height group=COMB2]
        Número de formas de elegir y ordenar \( k \) elementos de \( n \) sin repetir (permutaciones parciales).
        \[
            V_k^n = \frac{n!}{(n-k)!} = (n - k + 1) \cdot (n - k + 2) \cdots n
        \]
        Ejemplo: Números de 2 cifras de \{1,2,3\}: \( V_2^3 = 6 \).
    \end{tcolorbox}
    \begin{tcolorbox}[title=Variaciones con Repetición, equal height group=COMB2]
        Similar a variaciones, pero permitiendo repeticiones.
        \[
            VR_k^n = n^k
        \]
        Ejemplo: Números de 2 cifras de \{1,2,3\} con repetición: \( 3^2 = 9 \).
    \end{tcolorbox}
\end{multicols}

\begin{multicols}{2}
    \begin{tcolorbox}[title=Permutaciones sin Repetición, equal height group=COMB3]
        Número de formas de ordenar todos los \( n \) elementos de un conjunto. 
        \[
            P_n = n!
        \] \\
        Ejemplo: Permutaciones de \{1,2,3\}: \( 3! = 6 \).
    \end{tcolorbox}
    \begin{tcolorbox}[title=Permutaciones con Repetición, equal height group=COMB3]
        Número de permutaciones cuando hay elementos repetidos.
        \begin{align*}
            PR_{n_1, n_2, \dots, n_k}^n = \binom{n}{n_1 n_2 \dots n_k} = \frac{n!}{n_1! \cdot n_2! \cdots n_k!}, \\
            \quad \text{donde } n = n_1 + \cdots + n_k
        \end{align*} \\
        Ejemplo: Palabras con letras de ``PROBABILIDAD'' (12 letras, repeticiones específicas): \( \frac{12!}{{2!}^4 {1!}^4} = 29,937,600 \).
    \end{tcolorbox}
\end{multicols}

\begin{multicols}{2}
    \begin{tcolorbox}[title=Principio de la Suma, equal height group=COMB4]
        Suma los cardinales de conjuntos mutuamente exclusivos. \\
        Si \( A_1, \dots, A_n \) son disjuntos: \( \#(\cup_{i=1}^n A_i) = \sum_{i=1}^n \#(A_i) \).
        
    \end{tcolorbox}
    \begin{tcolorbox}[title=Principio del Producto, equal height group=COMB4]
        Corrige el solapamiento al contar uniones de conjuntos no disjuntos. \\
        \( \#(A_1 \cup A_2) = \#(A_1) + \#(A_2) - \#(A_1 \cap A_2) \).
    \end{tcolorbox}
\end{multicols}

\begin{tcolorbox}[title=Problema de las estrellas y barras]
Queremos encontrar el número de soluciones no negativas de $x_1 + x_2 + \cdots + x_k = n$, $x_i \geq 0$

Imaginamos que tenemos $n$ objetos idénticos que queremos repartir en $k$ cajones. Alineamos los $n$ objetos $\bigstar$ en fila:

\[
\bigstar \bigstar \bigstar \cdots \bigstar
\]

para separar estos objetos en $k$ cajones colocamos $k-1$ barras que marcan dónde termina cada grupo y empieza el siguiente.
Por ejemplo, $n = 7$ y $k = 3$:

\[
\bigstar \bigstar \bigstar \bigstar \mid \bigstar \bigstar \mid \bigstar
\]

Esto representa cada cajón $x_1 = 4, x_2 = 2, x_3 = 1$.

Cada solución corresponde a una secuencia de $n$ estrellas y $k-1$ barras en fila.

\begin{itemize}
    \item Número total de estrellas: $n$
    \item Número total de barras: $k-1$
    \item Número total de símbolos: $n + k - 1$
\end{itemize}

Así que el problema se reduce a: ¿de cuántas formas podemos elegir las posiciones de las $k-1$ barras entre $n+k-1$ lugares posibles?

\[
\binom{n+k-1}{k-1}
\]

\end{tcolorbox}

\sectiontitle{Sucesos}

\begin{tcolorbox}[title=Definiciones, equal height group=SUC1]
    \begin{tabular}{rl}
        \textbf{Experimento aleatorio}: & Experimento que, repetido en las mismas condiciones, puede dar resultados diferentes, \\
                                        & pero predecibles a largo plazo. \\
        \textbf{Suceso elemental}: & Cada posible resultado del experimento. \\
        \textbf{Espacio muestral \(\Omega\)}: & Conjunto de todos los sucesos elementales. \\
        \textbf{Suceso}: & Cualquier subconjunto de \(\Omega\). \\
        \textbf{Suceso seguro}: & \(\Omega\), \(P(\Omega) = 1\). \\
        \textbf{Suceso imposible}: & \(\emptyset\), \(P(\emptyset) = 0\).
    \end{tabular}
\end{tcolorbox}

\begin{multicols}{2}
    \begin{tcolorbox}[title=Operaciones, equal height group=SUC2]
        \begin{tabular}{rl}
            \textbf{Unión}: & \(A \cup B\) (si sucede \(A\) o \(B\)) \\
            \textbf{Intersección}: & \(A \cap B\) (si sucede \(A\) y \(B\)) \\
            \textbf{Complementario}: & \(A^c\) (si no sucede \(A\)) \\
            \textbf{Diferencia}: & \(A - B = A \cap B^c\) (sucede \(A\) y no \(B\)) \\
            \textbf{Incompatibles}: & \(A \cap B = \emptyset\)
        \end{tabular}
    \end{tcolorbox}
    \begin{tcolorbox}[title=Propiedades Operaciones, equal height group=SUC2]
        \begin{tabular}{rl}
            \textbf{Conmutativas}: & \(A \cup B = B \cup A\), \(A \cap B = B \cap A\) \\
            \textbf{Asociativas}: & \(A \cup (B \cup C) = (A \cup B) \cup C\) \\
                                  & \(A \cap (B \cap C) = (A \cap B) \cap C\) \\
            \textbf{Distributivas}: & \(A \cap (B \cup C) = (A \cap B) \cup (A \cap C)\) \\
                                    & \(A \cup (B \cap C) = (A \cup B) \cap (A \cup C)\) \\
            \textbf{Complementario}: & \({(A^c)}^c = A\) \\
            \textbf{De Morgan}: & \({(A \cup B)}^c = A^c \cap B^c\) \\
                                & \({(A \cap B)}^c = A^c \cup B^c\)
        \end{tabular}
    \end{tcolorbox}
\end{multicols}

\sectiontitle{Probabilidad}

\begin{multicols}{2}
    \begin{tcolorbox}[title=Definición, equal height group=PROB1]
        Para un espacio muestral finito \(\Omega\), una probabilidad \(P: \mathcal{P}(\Omega) \to [0,1]\) satisface:
        \begin{enumerate}[topsep=0pt]
            \item \(0 \leq P(A) \leq 1\) para todo suceso \(A\)
            \item \(P(\Omega) = 1\)
            \item Si \(A_1, \dots, A_n\) disjuntos, \\ \(P(A_1 \cup \cdots \cup A_n) = P(A_1) + \cdots + P(A_n)\)
        \end{enumerate}
        Si todos los elementales son equiprobables, \(P(A) = \frac{|A|}{|\Omega|}\).
    \end{tcolorbox}
    \begin{tcolorbox}[title=Propiedades, equal height group=PROB1]
        \begin{itemize}
            \item \(P(\emptyset) = 0\)
            \item \(P(A^c) = 1 - P(A)\)
            \item Si \(B \subseteq A\), entonces \(0 \leq P(B) \leq P(A)\)
            \item \(P(A - B) = P(A) - P(A \cap B)\)
            \item \(P(A \cup B) = P(A) + P(B) - P(A \cap B)\)
            \item Para tres conjuntos: \(P(A \cup B \cup C) = P(A) + P(B) + P(C) - P(A \cap B) - P(A \cap C) - P(B \cap C) + P(A \cap B \cap C)\)
        \end{itemize}
    \end{tcolorbox}
\end{multicols}

\begin{multicols}{2}
    \begin{tcolorbox}[title=Probabilidad Condicionada, equal height group=PROB2]
        Dado \(A, B\) con \(P(A) > 0\):
        \[
            P(B \mid A) = \frac{P(A \cap B)}{P(A)}.
        \]
        Propiedades (similar a probabilidad estándar):
        \begin{itemize}
            \item \(P(B^c \mid A) = 1 - P(B \mid A)\).
            \item \(P(B_1 \cup B_2 \mid A) = P(B_1 \mid A) + P(B_2 \mid A) - P(B_1 \cap B_2 \mid A)\).
        \end{itemize}
    \end{tcolorbox}
    \begin{tcolorbox}[title=Independencia, equal height group=PROB2]
        \(A\) y \(B\) son independientes si:
        \[
            P(A \cap B) = P(A) \cdot P(B)
        \]
        Equivalentes (si \(P(A), P(B) > 0\)):
        \begin{itemize}
            \item \(P(A \mid B) = P(A)\)
            \item \(P(B \mid A) = P(B)\)
        \end{itemize}
        Para familia \(\{A_1, \dots, A_n\}\): Para toda subfamilia, el producto de probabilidades individuales iguala la probabilidad de la intersección.
        También: Independencia se preserva con complementarios (e.g., \(A^c\) y \(B\) independientes si \(A\) y \(B\) lo son).
    \end{tcolorbox}
\end{multicols}

\begin{tcolorbox}[title=Probabilidad Total]
    Para \(A, B\):
    \[
        P(B) = P(B \cap A) + P(B \cap A^c) = P(A) \cdot P(B \mid A) + P(A^c) \cdot P(B \mid A^c).
    \]
    Para partición \(\{A_1, \dots, A_n\}\) de \(\Omega\):
    \[
        P(B) = \sum_{i=1}^n P(A_i) \cdot P(B \mid A_i).
    \]
\end{tcolorbox}

\begin{tcolorbox}[title=Teorema de Bayes]
    Para \(A, B\) con \(P(B) > 0\):
    \[
        P(A \mid B) = \frac{P(A) \cdot P(B \mid A)}{P(B)} = \frac{P(A) \cdot P(B \mid A)}{P(A) \cdot P(B \mid A) + P(A^c) \cdot P(B \mid A^c)}.
    \]
    Para partición \(\{A_1, \dots, A_n\}\):
    \[
        P(A_i \mid B) = \frac{P(A_i) \cdot P(B \mid A_i)}{\sum_{j=1}^n P(A_j) \cdot P(B \mid A_j)}.
    \]
\end{tcolorbox}




\sectiontitle{Variables Aleatorias}

\begin{tcolorbox}[title=Definiciones]
    \textbf{Variable aleatoria (v.a.):} Aplicación que asigna a cada elemento $\omega$ del espacio muestral $\Omega$ un número real $X(\omega) \in \mathbb{R}$.

    \begin{itemize}
        \item Se denota con letras mayúsculas: $X, Y, Z$.
        \item Sus valores se denotan en minúscula: $x, y, z$.
    \end{itemize}

    \textbf{Tipos de v.a.}
    \begin{itemize}
        \item \textbf{Discreta:} Toma valores numerables con probabilidad positiva.
        \item \textbf{Continua:} Sus valores están en intervalos y se describen con funciones de densidad.
    \end{itemize}
\end{tcolorbox}

\sectiontitle{Variables Aleatorias Discretas}

\begin{multicols}{2}
    \begin{tcolorbox}[title=Función de Probabilidad, equal height group=VA1]
        \[ P_X(x) = P(X = x)\]

        \textbf{Propiedades}:
        \begin{itemize}
            \item $\sum_{x \in D_X} P_X(x) = 1$
            \item $0 \leq P_X(x) \leq 1$
        \end{itemize}

        \textbf{Dominio}:
        \[ D_X(x) = \{ x \in \mathbb{R} \vert P_X(x) > 0 \} \]

        Ejemplo: dado justo
        \[
            P_X(x) = \begin{cases}
                \frac{1}{6}, & x = 1,2,3,4,5,6 \\
                0, & \text{en otro caso}
            \end{cases}
        \]

    \end{tcolorbox}
    \begin{tcolorbox}[title=Función de Distribución Acumulada (CDF), equal height group=VA1]
        \[ F_X(x) = P(X \leq x).  \]
        
        Ejemplo: dado
        \[
            F_X(x) = \begin{cases}
                0, & x < 1 \\
                \frac{k}{6}, & k \leq x < k+1, \; k=1,\dots,5 \\
                1, & x \geq 6
            \end{cases}
        \]

        Propiedades:
        \begin{itemize}
            \item $0 \leq F_X(x) \leq 1$ 
            \item $P(X > x) = 1 - P(X \leq x) = 1 - F_X(x)$ 
            \item Sean $a$ y $b$ tales que $a < b$: \\ $P(a < X \leq b) = P(X \leq b) - P(X \leq a) = F_X(b) - F_X(a)$
            \item $\lim\limits_{x\to-\infty}{F_X(x)} = 0$, $\lim\limits_{x\to\infty}{F_X(x)} = 1$
            \item Continua por la derecha: $\lim\limits_{x\to x_0^+}{F_X(x)} = F_X(x_0)$
        \end{itemize}
    \end{tcolorbox}
\end{multicols}

\begin{tcolorbox}[title=Valor Esperado, equal height group=VA2]
    El valor esperado, esperanza o media $E(X)$ de una v.a.\ discreta $X$ se define como:
    \[
        \overline{x} = \mu = \mu_X = E(X) = \sum\limits_{x \in X(\Omega)}{x \cdot P_X(x)}
    \]

    \textbf{Ejemplo}: lanzamos un dado $n$ veces y obtenemos unas frecuencias $n_i$ para $i = 1,\ldots,6$:
    \[
        \overline{x} = \frac{1 \cdot n_1 + 2 \cdot n_2 + 3 \cdot n_3 + 4 \cdot n_4 + 5 \cdot n_5 + 6 \cdot n_6}{n} = \sum\limits_{i = 1}^{6}{x \cdot P_X(x)}
    \]

    \textbf{Propiedades}:
    \begin{itemize}
        \item $E(k) = k$ para cualquier constante $k$
        \item Si $a \leq X \leq b$ entonces $a \leq E(X) \leq b$
        \item Si $X$ es una v.a.\ discreta que toma valores enteros no negativos entonces $E(X) = \sum\limits_{x=0}^{+\infty}{(1 - F_X(x))}$
        \item $E(g(X)) = \sum\limits_{x}{g(x) \cdot P_X(x)}$
    \end{itemize}
\end{tcolorbox}

\begin{tcolorbox}[title=Series Geométricas, equal height group=VA3]
    \begin{multicols}{2}
        \begin{tabular}{rl}
            \textbf{Progresión de razón $r$}: & $r^0, r^1, \ldots r^n$ \\
            \textbf{Serie}: & $\sum\limits_{k=0}^{+\infty}{r^k}$ \\
            \textbf{Sumas paricales}: & $\sum\limits_{k=n_0}^{n}{r^k} = \frac{r^{n_0} - r^n r}{1 - r}$ \\
            \textbf{Momento de orden $m$}: & $E({(X-C)}^m)$
        \end{tabular}
        \\
        \textbf{Propiedades}:
        \begin{itemize}
            \item Si se empieza en $n_0$ se tiene que $\sum\limits_{k=n_0}^{+\infty}{r^k} = {\frac{r^{n_0}}{1-r}}$
            \item Si $|r| < 1$ la serie es convergente y $\sum\limits_{k=0}^{+\infty}{r^k} = \frac{1}{1 - r}$ \\
            y sus derivadas también son convergentes:
            \begin{align*}
                {\left( \sum\limits_{k=0}^{+\infty}{r^k} \right)}'  = & \sum\limits_{k=1}^{+\infty}{kr^{k-1}}; {\left( \frac{1}{1-r} \right)}' = \frac{1}{{(1-r)}^2} \\
                {\left( \sum\limits_{k=0}^{+\infty}{r^k} \right)}'' = & \sum\limits_{k=2}^{+\infty}{k(k-1)r^{k-2}}; {\left( \frac{1}{1-r} \right)}'' = \frac{2}{{(1-r)}^3} \\
            \end{align*}
        \end{itemize}
    \end{multicols}
\end{tcolorbox}

\begin{tcolorbox}[title=Varianza, equal height group=VA4]
    \begin{multicols}{2}
        \textbf{Varianza}: $ \mathrm{Var}(X) = E({(X - E(X))}^2) = \sum\limits_{x}{{(x - E(X))}^2 P_X(x)} $ \\
        {\raggedleft\ $Var(X) = E(X^2) - (E(X))^2 = \sum\limits_x{x^2P_X(X) - (E(X))^2 }$ \par}
        \textbf{Desviación típica o estándar}: $ \sigma = \sigma_X = +\sqrt{Var(X)} $ \\
        \textbf{Propiedades}:
        \begin{itemize}
            \item $Var(X) \geq 0$
            \item $Var(\textit{cte}) = E(\textit{cte}^2) - ({E(\textit{cte})}^2) = \textit{cte}^2 - \textit{cte}^2 = 0$
            \item Si $Y = a + bX$:
            \begin{itemize}
                \item $E(Y) = E(a + bX) = a + bE(X) = a + b\mu_X$
                \item $Var(Y) = Var(a + bX) = b^2Var(x) = b^2\sigma_{X^-}^2$
                \item $\sigma_Y = \sqrt{Var(Y)} = \sqrt{b^2Var(X)} = |b|\sigma_{X^-}$
            \end{itemize}
        \end{itemize}
    \end{multicols}
    \vspace*{5pt}
\end{tcolorbox}

\sectiontitle{Variables Aleatorias Continuas}

\begin{tcolorbox}[title=Variables Aleatorias Continuas, equal height group=VAC1]
    \begin{multicols}{2}
        \textbf{Propiedades}
        \begin{itemize}
            \item $P_X(X = x_0) = F_X(x_0) - F_X(x_0^-) = 0$
            \item $P(X \leq b) = P(X < b)$
            \item $P(X < b) = P(X < a) + P(a < X < b)$
            \item $F_X(b) = F_X(a) + P(a < X < b)$
            \item $P(a < X < b) = F_X(b) - F_X(a) = P(a \leq X \leq b)$
        \end{itemize}
        \textbf{Función de densidad}
        \[
            \int_{-\infty}^{+\infty}{f_X(x)dx} = 1, f_X(x) \geq 0 \forall x \in \mathbb{R}
        \]
        \textbf{Función de distribución}
        \[
            F_X(x) = \int_{-\infty}^{x}f_X(t)dt, \forall x \in \mathbb{R}
        \]
    \end{multicols}
    \vspace*{1pt}
\end{tcolorbox}
\begin{tcolorbox}[title=Esperanza y Varianza, equal height group=VAC2]
    \begin{multicols}{2}
        \textbf{Esperanza}
        \[
            E(X) = \int_{-\infty}^{+\infty}{x \cdot f_X(x)dx}
        \]
        \textbf{Varianza}
        \[
            Var(X) = E((X - \mu_X)^2) = \int_{-\infty}^{+\infty}{(x - \mu_X)^2 \cdot f_X(x)dx}
        \]
        \textbf{Propiedades}
        \begin{itemize}
            \item $E(Y) = E(a + bX) = a + bE(X)$
            \item $Var(Y) = Var(a + bX) = b^2E(X)$
            \item $\sigma_Y = |b|\sigma_X$
            \item $Z = \frac{X - \mu_X}{\sigma_X} \implies E(Z) = 0, Var(Z) = 1$
        \end{itemize}
    \end{multicols}
    \vspace*{5pt}
\end{tcolorbox}

\sectiontitle{Distribuciones notables}

\begin{multicols}{2}
    \begin{tcolorbox}[title=Bernoulli, equal height group=DN1]
        Tirar una modeda y observar si sale cara o cruz (éxito o fracaso).
        \begin{tabular}{l}
            $X \sim Ber(p)$ \\
            $D_X = \{0,1\}$ \\
            $P_X(x) = P(X = x) = \begin{cases} 1-p=q & \text{si } x = 0 \\ p & \text{si } x = 1 \\ 0 & \text{en otro caso} \end{cases}$ \\
            $F_X(x) = P(X \leq x) = \begin{cases} 0 & \text{si } x < 0 \\ 1-p & \text{si } 0 \leq x < 1 \\ 1 & \text{si } x \leq 1\end{cases}$ \\
            $E(X) = p, Var(X) = p(1-p)$
        \end{tabular}
    \end{tcolorbox}
    \begin{tcolorbox}[title=Binomial, equal height group=DN1]
        Probabilidad de obtener $x$ caras al lanzar una moneda $n$ veces.
        \begin{tabular}{l}
            $X \sim B(n, p)$ \\
            $D_X = \{0,1,\ldots,n\}$ \\
            $P_X(x) = \begin{cases} \binom{n}{x}p^x(1-p)^{n-x} & \text{si } x \in [0, n] \\ 0 & \text{en otro caso} \end{cases}$ \\
            $F_X(x) = \begin{cases} 0 & \text{si } x < 0 \\ \sum\limits_{i=0}^{k}\binom{n}{i}p^i(1-p)^{n-i} & \text{si } k \leq x < k+1, \text{para } k \in [0, n] \\ 1 & \text{si } x \geq n \end{cases}$ \\
            $E(X) = np, Var(X) = np(1-p)$
        \end{tabular}
    \end{tcolorbox}
\end{multicols}

\pagebreak

\begin{tcolorbox}[title=Geométrica, equal height group=DN2]
    Tirar una moneda hasta que obtenemos la primera cara.
    \begin{multicols}{2}
        \begin{tabular}{l}
            $X \sim Ge(p)$ \\
            $X = \text{número de fracasos para conseguir primer éxito}$ \\
            $D_X = \{0,1,\ldots,n,\ldots\}$ \\
            $P_X(x) = P(X = x) = \begin{cases} (1-p)^xp & \text{si } x \geq 0 \\ 0 & \text{en otro caso} \end{cases}$ \\
            $F_X(x) = F(X \leq x) = \begin{cases} 0 & \text{si } x < 0 \\ 1-(1-p)^{k+1} & \text{si } \begin{cases} k \leq x < k+1 \\ \text{para } k \geq 0 \end{cases}\end{cases}$ \\
            $E(X) = \frac{1-p}{p}, Var(X) = \frac{1-p}{p^2}$
        \end{tabular}

        \columnbreak

        \begin{tabular}{l}
            $Y \sim Ge(p)$ \\
            $Y = \text{número de intentos para conseguir primer éxito}$ \\
            $D_Y = \{0,1,\ldots,n,\ldots\}$ \\
            $P_Y(y) = P(Y = y) = \begin{cases} (1-p)^{y-1} & \text{si } x \geq 1 \\ 0 & \text{en otro caso} \end{cases}$ \\
            $F_Y(y) = F(Y \leq y) = \begin{cases} 0 & \text{si } x < 1 \\ 1-(1-p)^k & \text{si } \begin{cases} k \leq x < k+1 \\ \text{para } k \geq 1 \end{cases}\end{cases}$ \\
            $E(Y) = \frac{1}{p}, Var(Y) = \frac{1-p}{p^2}$
        \end{tabular}
    \end{multicols}
    \vspace*{5pt}
\end{tcolorbox}

\begin{multicols}{2}
    \begin{tcolorbox}[title=Binomial Negativa, equal height group=DN3]
        Abrir dos puertas consecutivas con un conjunto de llaves sin memoria.
        \begin{tabular}{l}
            $X \sim BN(n, p)$ \\
            $X = \text{número de fracasos hasta conseguir el $n$-ésimo éxito}$ \\
            $D_X = \{0,1,2,3,\ldots\}$ \\
            $P_X(x) = P(X = x) = \begin{cases} \binom{x+n-1}{n-1} (1-p)^x p^n & \text{si } k \geq 0 \\ 0 & \text{en otro caso}\end{cases}$ \\
            $F_X(x) = P(X \leq x) = \begin{cases} \sum\limits_{i=0}^x P(X=i) & \text{si } \begin{cases} k \leq x < k+1, \\ k \geq 0 \end{cases}\end{cases}$ \\
            $E(X) = n \frac{1-p}{p}, Var(X) = n \frac{1-p}{p^2}$
        \end{tabular}
    \end{tcolorbox}
    \begin{tcolorbox}[title=Hipergeométrica, equal height group=DN3]
        Sacar bolas de una urna con reposición.
        \begin{tabular}{l}
            $X \sim H(m, n, k)$ \\
            $X = \begin{cases} \text{número de bolas blancas en $k$ extraccions} \\ \text{sín reposición de una urna con} \\ \text{$m$ bolas blancas y $n$ negras} \end{cases}$ \\
            $D_X = \{x \in \mathbb{N}\ \vert \max\{0,k-n\} \leq x \leq \min\{m, k\} \}$ \\
            $P_X(x) = \begin{cases} \frac{\binom{m}{x}\binom{n}{k-x}}{\binom{m+n}{k}} & \text{si } \max\{0, k-n\} \leq x \leq \min\{m, k\} \\ 0 & \text{en otro caso}\end{cases}$ \\
            $F_X(x) = P(X \leq x)$ \\
            $E(X) = \frac{km}{m + n}, Var(X) = k \frac{m}{m+n} (1-\frac{m}{m+n})\frac{m+n-k}{m+n-1}$
        \end{tabular}
    \end{tcolorbox}
\end{multicols}

\begin{tcolorbox}[title=Poisson, equal height group=DN4]
    \begin{multicols}{2}
        $X$ con distribución Poisson de media o promedio $\lambda$
        \begin{tabular}{l}
            $X \sim Po(\lambda)$ \\
            $D_X = \{0,1,\ldots\}$ \\
            $P_X(x) = P(X = x) = \begin{cases} \frac{\lambda^x}{x!}e^{-\lambda} & \text{si } x \geq 0 \\ 0 & \text{en otro caso}\end{cases}$ \\
            $F_X(x) = F(X \leq x) = \begin{cases} 0 & \text{si } x < 0 \\ \sum\limits_{i=0}^{x} P(X=i) & \text{si } \begin{cases} k \leq x < k + 1, \\ k \geq 0 \end{cases}\end{cases}$ \\
            $E(X) = \lambda, Var(X) = \lambda$
        \end{tabular}

        \columnbreak

        $X_t = $ número de eventos en el intervalo $(0, t] Po(\lambda\ t)$ donde $t$ promedio pro u.t.

        \begin{tabular}{l}
            $X \sim Po(\lambda\ t)$ \\
            $D_X = \{0,1,\ldots\}$ \\
            $P_X(x) = P(X = x) = \begin{cases} \frac{(\lambda\ t)^x}{x!}e^{-\lambda\ t} & \text{si } x \geq 0 \\ 0 & \text{en otro caso}\end{cases}$ \\
            $F_X(x) = F(X \leq x) = \begin{cases} 0 & \text{si } x < 0 \\ \sum\limits_{i=0}^{x} P(X=i) & \text{si } \begin{cases} k \leq x < k + 1, \\ k \geq 0 \end{cases}\end{cases}$ \\
            $E(X) = \lambda t, Var(X) = \lambda t$
        \end{tabular}
    \end{multicols}
    \vspace*{5pt}
\end{tcolorbox}

\begin{multicols}{2}
    \begin{tcolorbox}[title=Uniforme, equal height group=DN5]
        \begin{tabular}{l}
            $X \sim U(a, b)$ \\
            $D_X = (a, b)$ \\
            $f_X(x) = \begin{cases} \frac{1}{b-a} & \text{si } a < x < b \\ 0 & \text{en otro caso}\end{cases}$ \\
            $F_X(x) = P(X \leq x) = \begin{cases} 0 & \text{si } x \leq a \\ \frac{x-a}{b-a} & \text{si } a \leq x \leq b \\ 1 & \text{si } x \leq b\end{cases}$ \\
            $E(X) = \frac{a+b}{2}, Var(X) = \frac{(b-a)^2}{12}$
        \end{tabular}
    \end{tcolorbox}
    \begin{tcolorbox}[title=Exponencial, equal height group=DN5]
        \begin{tabular}{l}
            $X \sim Exp(\lambda)$ \\
            $D_X = (0, +\infty)$ \\
            $f_X(x) = \begin{cases} 0 & \text{si } x \leq 0 \\ \lambda e^{-\lambda x} & \text{si } x > 0\end{cases}$ \\
            $F_X(x) = P(X \leq x) = \begin{cases} 0 & \text{si } x \leq 0 \\ 1 - e^{-\lambda x} & \text{si } x > 0\end{cases}$ \\
            \\
            $E(X) = \frac{1}{\lambda}, Var(X) = \frac{1}{\lambda^2}$
        \end{tabular}
    \end{tcolorbox}
\end{multicols}

\begin{tcolorbox}[title=Normal o Gaussiana, equal height group=DN6]
    \begin{tabular}{l}
        $X \sim N(\mu, \sigma)$ \\
        $D_X = \mathbb{R} = (-\infty, +\infty)$ \\
        $f_X(x) = \frac{1}{\sqrt{2\pi \sigma^2}}e^{\frac{-(x-\mu)^2}{2\sigma^2}}, \forall x \in \mathbb{R}$ \\
        $F_X(x) = P(X \leq x) = \int_{-\infty}^{x} f_X(t)dt$ \\
        \\
        $E(X) = \mu, Var(X) = \sigma^2$
    \end{tabular}
\end{tcolorbox}

\pagebreak

\sectiontitle{Variables Aleatorias Multidimensionales}

\begin{tcolorbox}[title=Variables Aleatorias Bidimensionales]
    \begin{multicols}{2}
        \textbf{Definiciones}
        \begin{itemize}
            \item \textbf{Discreta}: su conjunto de valores es $\mathbb{R}^2, (X, Y)(\Omega)$ es un conjunto finito o numerable.
            \item \textbf{Continua}: su conjunto de valores es $\mathbb{R}^2, (X, Y)(\Omega)$ es un conjunto es un producto de intervalos.
            \item \textbf{Heterogénea}: $X$ e $Y$ no comportan ser continuas o discretas.
        \end{itemize}
    
        \textbf{Dominio} \\
        $D_{XY} = \{(x, y) \in \mathbb{R}^2 \mid P_{XY}(x, y) > 0\}$
    
        \textbf{Función de probabilidad} \\
        $P_{XY}(x, y) = P(X = x, Y = y)$ \\
        $\sum\limits_i\sum\limits_j P_{XY}(x_i, y_j) = 1$ \\
        $P((X, Y) \in B) = \sum\limits_{(x_i, x_j) \in B} P_{XY}(x_i, y_j), B \subset D_{XY}$

        \columnbreak
    
        \textbf{Función de distribución acumulada} \\
        $F_{XY} = P(X \leq x, Y \leq y) = \sum\limits_{x_i \leq x, y_j \leq y} P_{XY}(x_i, y_j)$
    
        \textbf{Función de densidad} \\
        $f_{XY} : \mathbb{R} \times \mathbb{R} \to [0, +\infty)$ \\
        \begin{itemize}
            \item $f_{XY}(x, y) > 0, \forall (x, y) \in D_{XY}$
            \item $\int_{-\infty}^{+\infty} \int_{-\infty}^{+\infty} f_{XY}(t_x, t_y)dt_x dt_y = 1$
            \item $F_{XY}(x, y) = \int_{-\infty}^{x} \int_{-\infty}^{y} f_{XY}(t_x, t_y)dt_x dt_y$
        \end{itemize}
    \end{multicols}
\end{tcolorbox}

\begin{tcolorbox}[title=Distribuciones Marginales]
    \begin{multicols}{2}
        \textbf{Probabilidad Marginal} \\
        $P_X(x_i) = \sum\limits_j P_{XY}(x_i, y_j)$ \\
        $P_Y(y_j) = \sum\limits_i P_{XY}(x_i, y_j)$

        \textbf{Densidad Marginal} \\
        $f_X(x) = \int_{-\infty}^{+\infty} f_XY{x, y} dy$ \\
        $f_Y(y) = \int_{-\infty}^{+\infty} f_XY{x, y} dx$
    \end{multicols}

    \hrule

    \begin{multicols}{2}
        {\large\bfseries\ Discretas}\\
        \textbf{Independencia} \\
        $X$ e $Y$ son independientes si se cumple alguna de las condiciones:
        \begin{itemize}
            \item $P_{XY}(x_i, y_j) = P_X(x_i) \cdot P_Y{y_j}$
            \item $P(X = x_i, Y = y_j) = P(X = x_i) \cdot P{Y = y_j}$
            \item $F_{XY}(x, y) = F_X(x) \cdot F_Y(y)$
        \end{itemize}

        \textbf{Esperanza y varianza} \\
        \begin{itemize}
            \item $E(X) = \sum\limits_{x \in D_X} x \cdot P_X(x)$
            \item $E(Y) = \sum\limits_{y \in D_Y} y \cdot P_Y(y)$
            \item $\sigma_X^2 = Var(X) = E(X - E(X)) = E(X) - E(X)^2$
            \item $\sigma_Y^2 = Var(Y) = E(Y - E(Y)) = E(Y) - E(Y)^2$
        \end{itemize}

        \textbf{Distribuciones condicionales} \\
        Dado un valor fijo $y \in D_Y$:
        \[
            P(X = x \mid Y = y) = \frac{P_{XY}(x, y)}{P_Y(y)}
        \]
        Dado un valor fijo $x \in D_X$:
        \[
            P(Y = y \mid X = x) = \frac{P_{XY}(x, y)}{P_X(x)}
        \]
        Independientes si y sólo si:
        \begin{itemize}
            \item $P(X = x \mid Y = y) = P(X = x)$
            \item $P(Y = y \mid X = x) = P(Y = y)$
        \end{itemize}
        Esperanza:
        \[
            E(X \mid Y=y) = \sum\limits_{x \in D_X} x \cdot P(X = x \mid Y = y)
        \]

        \columnbreak

        {\large\bfseries\ Continuas}\\
        \textbf{Independencia} \\
        $X$ e $Y$ son independientes si se cumple alguna de las condiciones:
        \begin{itemize}
            \item $f_{XY}(x, y) = f_X(x) \cdot f_Y(y), \forall (x, y) \in D_{XY}$
            \item $F_{XY}(x, y) = F_X(x) \cdot F_Y(y), \forall (x, y) \in D_{XY}$
        \end{itemize}

        \textbf{Esperanza y varianza} \\
        \begin{itemize}
            \item $E(X) = \int_{-\infty}^{+\infty} x \cdot f_X(x)dx$
            \item $E(Y) = \int_{-\infty}^{+\infty} y \cdot f_Y(y)dy$
            \item $\sigma_X^2 = Var(X) = E((X - E(X))^2) = E(X^2) - E(X)^2$
            \item $\sigma_Y^2 = Var(Y) = E((Y - E(Y))^2) = E(Y^2) - E(Y)^2$
        \end{itemize}

        \textbf{Distribuciones condicionales} \\
        Dado un valor fijo $y \in D_Y$:
        \[
            f_{X \mid Y=y}(x) = \frac{f_{XY}(x, y)}{f_Y(y)}
        \]
        Dado un valor fijo $x \in D_X$:
        \[
            f_{Y \mid X=x}(y) = \frac{f_{XY}(x, y)}{f_X(x)}
        \]
        Independientes si y sólo si:
        \begin{itemize}
            \item $f_{X \mid Y=y}(x) = f_X(x)$
            \item $f_{Y \mid X=x}(y) = f_Y(y)$
        \end{itemize}
        Esperanza:
        \[
            E(X \mid Y=y) = \int_{-\infty}^{+\infty} x \cdot f_{X \mid Y=y}(x)dx
        \]
    \end{multicols}
    $X$ e $Y$ independientes si:
    \begin{itemize}
        \item $E(X \mid Y = y) = E(X)$
        \item $E(Y \mid X = x) = E(Y)$
    \end{itemize}
\end{tcolorbox}

\pagebreak

\sectiontitle{Estimación Puntual}

\begin{tcolorbox}[title=Estadística Inferencial]
    \textbf{Problema}: Queremos conocer una característica de una población, pero no podemos medirla en todos los individuos. Extraemos una muestra aleatoria e inferimos el valor para toda la población.
\end{tcolorbox}

\begin{multicols}{2}
    \begin{tcolorbox}[title=Muestra Aleatoria Simple (m.a.s.), equal height group=EST1]
        \textbf{Definición intuitiva}: De una población de $N$ individuos, repetimos $n$ veces el proceso de escoger equiprobablemente un individuo (con reposición).
        
        \textbf{Definición formal}: Un conjunto de $n$ variables aleatorias independientes $X_1, \ldots, X_n$, todas con la distribución de $X$.
        
        \textbf{Realización}: Los $n$ valores $x_1, \ldots, x_n$ que toman las v.a. $X_1, \ldots, X_n$.
    \end{tcolorbox}
    \begin{tcolorbox}[title=Estadístico (Estimador Puntual), equal height group=EST1]
        \textbf{Definición}: Una función aplicada a la muestra $X_1, \ldots, X_n$:
        \[
            T = f(X_1, \ldots, X_n)
        \]
        
        \textbf{Propiedades}:
        \begin{itemize}
            \item Es una variable aleatoria con distribución, esperanza, etc.
            \item La distribución muestral de $T$ es la distribución de esta v.a.
            \item Error estándar de $T$: desviación típica de $T$ ($\sigma_T$)
        \end{itemize}
        
        \textbf{Convenio}: Estadísticos en mayúsculas ($\bar{X}$), realizaciones en minúsculas ($\bar{x}$).
    \end{tcolorbox}
\end{multicols}

\begin{tcolorbox}[title=Media Muestral]
    \textbf{Definición}: Sea $X_1, \ldots, X_n$ una m.a.s. de tamaño $n$ de una v.a. $X$ con esperanza $\mu_X$ y desviación típica $\sigma_X$:
    \[
        \bar{X} = \frac{X_1 + \cdots + X_n}{n}
    \]
    
    \textbf{Propiedades}:
    \begin{itemize}
        \item $E(\bar{X}) = \mu_X$ (estimador insesgado)
        \item $\sigma_{\bar{X}} = \frac{\sigma_X}{\sqrt{n}}$ (error estándar)
        \item $\bar{X}$ es un estimador puntual de $\mu_X$
    \end{itemize}
    
    \textbf{Interpretación}: Si tomamos muchas veces una m.a.s. y calculamos la media muestral, el promedio de estas medias será aproximadamente $\mu_X$.
    
    \textbf{Teorema Central del Límite}: Si $n$ es suficientemente grande, $\bar{X}$ sigue aproximadamente una distribución normal:
    \[
        \bar{X} \approx N\left(\mu_X, \frac{\sigma_X}{\sqrt{n}}\right)
    \]
\end{tcolorbox}

\begin{multicols}{2}
    \begin{tcolorbox}[title=Propiedades de la Media Muestral, equal height group=EST2]
        \textbf{Insesgadez}: $E(\bar{X}) = \mu_X$
        
        \textbf{Consistencia}: A medida que $n$ aumenta, $\sigma_{\bar{X}} = \frac{\sigma_X}{\sqrt{n}}$ disminuye.
        
        \textbf{Eficiencia}: Entre los estimadores insesgados, la media muestral tiene varianza mínima (bajo ciertas condiciones).
    \end{tcolorbox}
    \begin{tcolorbox}[title=Notación y Convenios, equal height group=EST2]
        \textbf{Estadísticos} (v.a.): $X_1, \ldots, X_n$, $\bar{X}$, $S^2$
        
        \textbf{Realizaciones} (valores): $x_1, \ldots, x_n$, $\bar{x}$, $s^2$
        
        \textbf{Muestra sin reposición}: Si $N \gg n$, los resultados son aproximadamente los mismos que para m.a.s.
    \end{tcolorbox}
\end{multicols}

\pagebreak

\sectiontitle{Intervalos de Confianza}

\begin{tcolorbox}[title=Definición de Intervalo de Confianza]
    \textbf{Concepto}: Un intervalo de confianza $(1-\alpha)$ para un parámetro $\theta$ es un intervalo aleatorio $[L, U]$ tal que:
    \[
        P(L \leq \theta \leq U) = 1 - \alpha
    \]
    
    \textbf{Interpretación}: Si construimos muchos intervalos de confianza $(1-\alpha)$ a partir de diferentes muestras, aproximadamente el $(1-\alpha) \times 100\%$ de ellos contendrá el verdadero valor del parámetro $\theta$.
    
    \textbf{Nivel de confianza}: $1-\alpha$ (típicamente 0.95, 0.99)
    
    \textbf{Nivel de significación}: $\alpha$ (típicamente 0.05, 0.01)
\end{tcolorbox}

\begin{multicols}{2}
    \begin{tcolorbox}[title=Construcción de I.C., equal height group=IC1]
        \textbf{Pasos}:
        \begin{enumerate}
            \item Elegir nivel de confianza $1-\alpha$
            \item Encontrar estadístico pivote
            \item Determinar distribución del pivote
            \item Calcular cuantiles $z_{\alpha/2}$ o $t_{\alpha/2}$
            \item Construir intervalo
        \end{enumerate}
        
        \textbf{Forma general}:
        \[
            \text{Estimador} \pm \text{Margen de error}
        \]
    \end{tcolorbox}
    \begin{tcolorbox}[title=Cuantiles Importantes, equal height group=IC1]
        \textbf{Distribución Normal}:
        \begin{itemize}
            \item $z_{0.025} = 1.96$ (95\%)
            \item $z_{0.005} = 2.576$ (99\%)
            \item $z_{0.05} = 1.645$ (90\%)
        \end{itemize}
        
        \textbf{Distribución t de Student}:
        \begin{itemize}
            \item Depende de grados de libertad $\nu = n-1$
            \item $t_{\alpha/2, \nu}$ se obtiene de tablas
        \end{itemize}
    \end{tcolorbox}
\end{multicols}

\begin{tcolorbox}[title=Intervalo de Confianza para la Media $\mu$]
    \textbf{Caso 1: Varianza conocida $\sigma^2$} (o $n$ grande):
    \[
        \bar{X} \pm z_{\alpha/2} \cdot \frac{\sigma}{\sqrt{n}}
    \]
    Intervalo: $\left[\bar{x} - z_{\alpha/2} \frac{\sigma}{\sqrt{n}}, \bar{x} + z_{\alpha/2} \frac{\sigma}{\sqrt{n}}\right]$
    
    \textbf{Caso 2: Varianza desconocida, muestra pequeña} ($n < 30$):
    \[
        \bar{X} \pm t_{\alpha/2, n-1} \cdot \frac{S}{\sqrt{n}}
    \]
    donde $S = \sqrt{\frac{1}{n-1}\sum_{i=1}^n (X_i - \bar{X})^2}$ es la desviación típica muestral.
    
    Intervalo: $\left[\bar{x} - t_{\alpha/2, n-1} \frac{s}{\sqrt{n}}, \bar{x} + t_{\alpha/2, n-1} \frac{s}{\sqrt{n}}\right]$
    
    \textbf{Amplitud del intervalo}: $2 \cdot z_{\alpha/2} \frac{\sigma}{\sqrt{n}}$ (o con $t$ y $s$)
\end{tcolorbox}

\begin{multicols}{2}
    \begin{tcolorbox}[title=I.C. para Proporción $p$, equal height group=IC2]
        \textbf{Condición}: $n$ grande ($np \geq 5$ y $n(1-p) \geq 5$)
        
        \textbf{Estimador}: $\hat{p} = \frac{X}{n}$ donde $X \sim B(n, p)$
        
        \textbf{Intervalo}:
        \[
            \hat{p} \pm z_{\alpha/2} \sqrt{\frac{\hat{p}(1-\hat{p})}{n}}
        \]
        
        \textbf{Forma}: $\left[\hat{p} - z_{\alpha/2} \sqrt{\frac{\hat{p}(1-\hat{p})}{n}}, \hat{p} + z_{\alpha/2} \sqrt{\frac{\hat{p}(1-\hat{p})}{n}}\right]$
    \end{tcolorbox}
    \begin{tcolorbox}[title=I.C. para Varianza $\sigma^2$, equal height group=IC2]
        \textbf{Condición}: Población normal
        
        \textbf{Estadístico}: $\frac{(n-1)S^2}{\sigma^2} \sim \chi^2_{n-1}$
        
        \textbf{Intervalo}:
        \[
            \left[\frac{(n-1)s^2}{\chi^2_{\alpha/2, n-1}}, \frac{(n-1)s^2}{\chi^2_{1-\alpha/2, n-1}}\right]
        \]
        
        \textbf{Nota}: $\chi^2_{\alpha/2, n-1}$ y $\chi^2_{1-\alpha/2, n-1}$ son cuantiles de chi-cuadrado con $n-1$ grados de libertad.
    \end{tcolorbox}
\end{multicols}

\begin{tcolorbox}[title=Distribución t de Student]
    \textbf{Definición}: Si $Z \sim N(0,1)$ y $V \sim \chi^2_\nu$ son independientes, entonces:
    \[
        T = \frac{Z}{\sqrt{V/\nu}} \sim t_\nu
    \]
    
    \textbf{Propiedades}:
    \begin{itemize}
        \item Simétrica alrededor de 0
        \item $E(T) = 0$ (si $\nu > 1$)
        \item $Var(T) = \frac{\nu}{\nu-2}$ (si $\nu > 2$)
        \item Cuando $\nu \to \infty$, $t_\nu \to N(0,1)$
    \end{itemize}
    
    \textbf{Uso}: Cuando la varianza poblacional es desconocida y la muestra es pequeña.
\end{tcolorbox}

\begin{multicols}{2}
    \begin{tcolorbox}[title=Tamaño Muestral, equal height group=IC3]
        \textbf{Para media} (varianza conocida):
        \[
            n = \left(\frac{z_{\alpha/2} \cdot \sigma}{E}\right)^2
        \]
        donde $E$ es el error máximo deseado.
        
        \textbf{Para proporción}:
        \[
            n = \frac{z_{\alpha/2}^2 \cdot p(1-p)}{E^2}
        \]
        Si $p$ es desconocido, usar $p = 0.5$ (máxima varianza).
    \end{tcolorbox}
    \begin{tcolorbox}[title=Interpretación y Consideraciones, equal height group=IC3]
        \textbf{Interpretación correcta}: "Tenemos un 95\% de confianza de que el intervalo contiene el verdadero parámetro."
        
        \textbf{Interpretación incorrecta}: "El parámetro tiene un 95\% de probabilidad de estar en el intervalo."
        
        \textbf{Factores que afectan la amplitud}:
        \begin{itemize}
            \item Nivel de confianza: mayor confianza $\Rightarrow$ mayor amplitud
            \item Tamaño muestral: mayor $n$ $\Rightarrow$ menor amplitud
            \item Varianza: mayor varianza $\Rightarrow$ mayor amplitud
        \end{itemize}
    \end{tcolorbox}
\end{multicols}

\pagebreak

\sectiontitle{Contraste de Hipótesis}

\begin{tcolorbox}[title=Definiciones Básicas]
    \textbf{Hipótesis estadística}: Afirmación sobre un parámetro poblacional que se desea verificar.
    
    \textbf{Hipótesis nula ($H_0$)}: Afirmación que se asume verdadera inicialmente. Representa "no hay efecto" o "no hay diferencia".
    
    \textbf{Hipótesis alternativa ($H_1$ o $H_a$)}: Afirmación que contradice $H_0$. Representa el efecto o diferencia que se desea detectar.
    
    \textbf{Estadístico de prueba}: Función de la muestra que se usa para tomar la decisión.
    
    \textbf{Región crítica (o de rechazo)}: Conjunto de valores del estadístico que llevan a rechazar $H_0$.
    
    \textbf{Región de aceptación}: Conjunto de valores del estadístico que llevan a no rechazar $H_0$.
\end{tcolorbox}

\begin{multicols}{2}
    \begin{tcolorbox}[title=Tipos de Contrastes, equal height group=CH1]
        \textbf{Unilateral (cola derecha)}:
        \begin{itemize}
            \item $H_0: \theta = \theta_0$
            \item $H_1: \theta > \theta_0$
        \end{itemize}
        
        \textbf{Unilateral (cola izquierda)}:
        \begin{itemize}
            \item $H_0: \theta = \theta_0$
            \item $H_1: \theta < \theta_0$
        \end{itemize}
        
        \textbf{Bilateral}:
        \begin{itemize}
            \item $H_0: \theta = \theta_0$
            \item $H_1: \theta \neq \theta_0$
        \end{itemize}
    \end{tcolorbox}
    \begin{tcolorbox}[title=Errores en el Contraste, equal height group=CH1]
        \textbf{Error Tipo I ($\alpha$)}: Rechazar $H_0$ cuando es verdadera.
        \[
            \alpha = P(\text{Rechazar } H_0 \mid H_0 \text{ es verdadera})
        \]
        
        \textbf{Error Tipo II ($\beta$)}: No rechazar $H_0$ cuando es falsa.
        \[
            \beta = P(\text{No rechazar } H_0 \mid H_0 \text{ es falsa})
        \]
        
        \textbf{Nivel de significación}: Probabilidad máxima permitida de cometer Error Tipo I. Valores comunes: $\alpha = 0.05, 0.01, 0.10$.
    \end{tcolorbox}
\end{multicols}

\begin{multicols}{2}
    \begin{tcolorbox}[title=Potencia del Test, equal height group=CH2]
        \textbf{Definición}: Probabilidad de rechazar correctamente $H_0$ cuando $H_1$ es verdadera.
        \[
            \text{Potencia} = 1 - \beta = P(\text{Rechazar } H_0 \mid H_1 \text{ es verdadera})
        \]
        
        \textbf{Factores que aumentan la potencia}:
        \begin{itemize}
            \item Mayor tamaño muestral $n$
            \item Mayor diferencia entre $\theta_0$ y $\theta_1$
            \item Mayor nivel de significación $\alpha$
            \item Menor varianza poblacional
        \end{itemize}
    \end{tcolorbox}
    \begin{tcolorbox}[title=p-valor, equal height group=CH2]
        \textbf{Definición}: Probabilidad de obtener un resultado tan extremo como el observado (o más extremo), suponiendo que $H_0$ es verdadera.
        
        \textbf{Regla de decisión}:
        \begin{itemize}
            \item Si $p\text{-valor} < \alpha$: Rechazar $H_0$
            \item Si $p\text{-valor} \geq \alpha$: No rechazar $H_0$
        \end{itemize}
        
        \textbf{Interpretación}:
        \begin{itemize}
            \item $p\text{-valor}$ pequeño: Evidencia fuerte contra $H_0$
            \item $p\text{-valor}$ grande: Poca evidencia contra $H_0$
        \end{itemize}
    \end{tcolorbox}
\end{multicols}

\begin{tcolorbox}[title=Procedimiento General para Contrastes]
    \textbf{Pasos}:
    \begin{enumerate}
        \item Formular $H_0$ y $H_1$
        \item Elegir el nivel de significación $\alpha$
        \item Seleccionar el estadístico de prueba adecuado
        \item Determinar la región crítica (o calcular el $p$-valor)
        \item Calcular el valor del estadístico con los datos muestrales
        \item Tomar decisión: rechazar o no rechazar $H_0$
        \item Interpretar los resultados en el contexto del problema
    \end{enumerate}
\end{tcolorbox}

\sectiontitle{Contrastes para la Media Poblacional $\mu$}

\begin{multicols}{2}
    \begin{tcolorbox}[title=Caso 1: Varianza conocida $\sigma^2$, equal height group=CH3]
        \textbf{Condiciones}: Población normal o $n$ grande ($n \geq 30$)
        
        \textbf{Estadístico de prueba}:
        \[
            Z = \frac{\bar{X} - \mu_0}{\sigma/\sqrt{n}} \sim N(0,1)
        \]
        
        \textbf{Regiones críticas}:
        \begin{itemize}
            \item $H_1: \mu > \mu_0$: Rechazar si $Z > z_\alpha$
            \item $H_1: \mu < \mu_0$: Rechazar si $Z < -z_\alpha$
            \item $H_1: \mu \neq \mu_0$: Rechazar si $|Z| > z_{\alpha/2}$
        \end{itemize}
        
        \textbf{p-valor}:
        \begin{itemize}
            \item $H_1: \mu > \mu_0$: $p = P(Z > z_{\text{obs}})$
            \item $H_1: \mu < \mu_0$: $p = P(Z < z_{\text{obs}})$
            \item $H_1: \mu \neq \mu_0$: $p = 2P(Z > |z_{\text{obs}}|)$
        \end{itemize}
    \end{tcolorbox}
    \begin{tcolorbox}[title=Caso 2: Varianza desconocida, equal height group=CH3]
        \textbf{Condiciones}: Población normal, muestra pequeña ($n < 30$)
        
        \textbf{Estadístico de prueba}:
        \[
            T = \frac{\bar{X} - \mu_0}{S/\sqrt{n}} \sim t_{n-1}
        \]
        donde $S = \sqrt{\frac{1}{n-1}\sum_{i=1}^n (X_i - \bar{X})^2}$
        
        \textbf{Regiones críticas}:
        \begin{itemize}
            \item $H_1: \mu > \mu_0$: Rechazar si $T > t_{\alpha, n-1}$
            \item $H_1: \mu < \mu_0$: Rechazar si $T < -t_{\alpha, n-1}$
            \item $H_1: \mu \neq \mu_0$: Rechazar si $|T| > t_{\alpha/2, n-1}$
        \end{itemize}
        
        \textbf{Nota}: Si $n$ es grande ($n \geq 30$), se puede usar $Z$ aunque $\sigma$ sea desconocida.
    \end{tcolorbox}
\end{multicols}

\sectiontitle{Contrastes para la Proporción Poblacional $p$}

\begin{tcolorbox}[title=Contraste para una Proporción]
    \textbf{Condiciones}: $n$ grande ($np_0 \geq 5$ y $n(1-p_0) \geq 5$)
    
    \textbf{Estadístico de prueba}:
    \[
        Z = \frac{\hat{p} - p_0}{\sqrt{\frac{p_0(1-p_0)}{n}}} \sim N(0,1)
    \]
    donde $\hat{p} = \frac{X}{n}$ es la proporción muestral.
    
    \textbf{Regiones críticas}:
    \begin{itemize}
        \item $H_1: p > p_0$: Rechazar si $Z > z_\alpha$
        \item $H_1: p < p_0$: Rechazar si $Z < -z_\alpha$
        \item $H_1: p \neq p_0$: Rechazar si $|Z| > z_{\alpha/2}$
    \end{itemize}
    
    \textbf{p-valor}:
    \begin{itemize}
        \item $H_1: p > p_0$: $p = P(Z > z_{\text{obs}})$
        \item $H_1: p < p_0$: $p = P(Z < z_{\text{obs}})$
        \item $H_1: p \neq p_0$: $p = 2P(Z > |z_{\text{obs}}|)$
    \end{itemize}
\end{tcolorbox}

\sectiontitle{Contrastes para la Varianza Poblacional $\sigma^2$}

\begin{tcolorbox}[title=Contraste para la Varianza]
    \textbf{Condiciones}: Población normal
    
    \textbf{Estadístico de prueba}:
    \[
        \chi^2 = \frac{(n-1)S^2}{\sigma_0^2} \sim \chi^2_{n-1}
    \]
    donde $S^2 = \frac{1}{n-1}\sum_{i=1}^n (X_i - \bar{X})^2$
    
    \textbf{Regiones críticas}:
    \begin{itemize}
        \item $H_1: \sigma^2 > \sigma_0^2$: Rechazar si $\chi^2 > \chi^2_{\alpha, n-1}$
        \item $H_1: \sigma^2 < \sigma_0^2$: Rechazar si $\chi^2 < \chi^2_{1-\alpha, n-1}$
        \item $H_1: \sigma^2 \neq \sigma_0^2$: Rechazar si $\chi^2 < \chi^2_{1-\alpha/2, n-1}$ o $\chi^2 > \chi^2_{\alpha/2, n-1}$
    \end{itemize}
\end{tcolorbox}

\sectiontitle{Contrastes para Dos Muestras}

\begin{multicols}{2}
    \begin{tcolorbox}[title=Diferencia de Medias (Muestras Independientes), equal height group=CH4]
        \textbf{Caso: Varianzas conocidas} \\
        Estadístico:
        \[
            Z = \frac{\bar{X}_1 - \bar{X}_2 - (\mu_1 - \mu_2)_0}{\sqrt{\frac{\sigma_1^2}{n_1} + \frac{\sigma_2^2}{n_2}}} \sim N(0,1)
        \]
        
        \textbf{Caso: Varianzas desconocidas pero iguales} \\
        Estadístico:
        \[
            T = \frac{\bar{X}_1 - \bar{X}_2 - (\mu_1 - \mu_2)_0}{S_p\sqrt{\frac{1}{n_1} + \frac{1}{n_2}}} \sim t_{n_1+n_2-2}
        \]
        donde $S_p^2 = \frac{(n_1-1)S_1^2 + (n_2-1)S_2^2}{n_1+n_2-2}$ (varianza combinada)
        
        \textbf{Caso: Varianzas desconocidas y diferentes} \\
        Estadístico (aproximado):
        \[
            T = \frac{\bar{X}_1 - \bar{X}_2 - (\mu_1 - \mu_2)_0}{\sqrt{\frac{S_1^2}{n_1} + \frac{S_2^2}{n_2}}} \sim t_\nu
        \]
        donde $\nu = \frac{\left(\frac{S_1^2}{n_1} + \frac{S_2^2}{n_2}\right)^2}{\frac{(S_1^2/n_1)^2}{n_1-1} + \frac{(S_2^2/n_2)^2}{n_2-1}}$ (Welch)
    \end{tcolorbox}
    \begin{tcolorbox}[title=Diferencia de Medias (Muestras Pareadas), equal height group=CH4]
        \textbf{Condiciones}: Población normal, muestras dependientes
        
        \textbf{Definición}: $D_i = X_{1i} - X_{2i}$ (diferencias pareadas)
        
        \textbf{Estadístico de prueba}:
        \[
            T = \frac{\bar{D} - \mu_{D_0}}{S_D/\sqrt{n}} \sim t_{n-1}
        \]
        donde:
        \begin{itemize}
            \item $\bar{D} = \frac{1}{n}\sum_{i=1}^n D_i$ (media de diferencias)
            \item $S_D = \sqrt{\frac{1}{n-1}\sum_{i=1}^n (D_i - \bar{D})^2}$ (desviación típica de diferencias)
        \end{itemize}
        
        \textbf{Hipótesis típicas}:
        \begin{itemize}
            \item $H_0: \mu_D = 0$ (no hay diferencia)
            \item $H_1: \mu_D \neq 0$ (hay diferencia)
        \end{itemize}
    \end{tcolorbox}
\end{multicols}

\begin{multicols}{2}
    \begin{tcolorbox}[title=Diferencia de Proporciones, equal height group=CH5]
        \textbf{Condiciones}: $n_1$ y $n_2$ grandes
        
        \textbf{Estadístico de prueba}:
        \[
            Z = \frac{\hat{p}_1 - \hat{p}_2 - (p_1 - p_2)_0}{\sqrt{\hat{p}(1-\hat{p})\left(\frac{1}{n_1} + \frac{1}{n_2}\right)}} \sim N(0,1)
        \]
        donde $\hat{p} = \frac{X_1 + X_2}{n_1 + n_2}$ (proporción combinada)
        
        \textbf{Regiones críticas}:
        \begin{itemize}
            \item $H_1: p_1 > p_2$: Rechazar si $Z > z_\alpha$
            \item $H_1: p_1 < p_2$: Rechazar si $Z < -z_\alpha$
            \item $H_1: p_1 \neq p_2$: Rechazar si $|Z| > z_{\alpha/2}$
        \end{itemize}
    \end{tcolorbox}
    \begin{tcolorbox}[title=Cociente de Varianzas, equal height group=CH5]
        \textbf{Condiciones}: Poblaciones normales e independientes
        
        \textbf{Estadístico de prueba}:
        \[
            F = \frac{S_1^2/\sigma_1^2}{S_2^2/\sigma_2^2} = \frac{S_1^2}{S_2^2} \sim F_{n_1-1, n_2-1}
        \]
        (bajo $H_0: \sigma_1^2 = \sigma_2^2$)
        
        \textbf{Regiones críticas}:
        \begin{itemize}
            \item $H_1: \sigma_1^2 > \sigma_2^2$: Rechazar si $F > F_{\alpha, n_1-1, n_2-1}$
            \item $H_1: \sigma_1^2 < \sigma_2^2$: Rechazar si $F < F_{1-\alpha, n_1-1, n_2-1}$
            \item $H_1: \sigma_1^2 \neq \sigma_2^2$: Rechazar si $F < F_{1-\alpha/2, n_1-1, n_2-1}$ o $F > F_{\alpha/2, n_1-1, n_2-1}$
        \end{itemize}
        
        \textbf{Nota}: $F_{1-\alpha, \nu_1, \nu_2} = \frac{1}{F_{\alpha, \nu_2, \nu_1}}$
    \end{tcolorbox}
\end{multicols}

\begin{tcolorbox}[title=Distribución F de Fisher]
    \textbf{Definición}: Si $U \sim \chi^2_{\nu_1}$ y $V \sim \chi^2_{\nu_2}$ son independientes, entonces:
    \[
        F = \frac{U/\nu_1}{V/\nu_2} \sim F_{\nu_1, \nu_2}
    \]
    
    \textbf{Propiedades}:
    \begin{itemize}
        \item $F_{\nu_1, \nu_2}$ tiene dos parámetros: grados de libertad del numerador ($\nu_1$) y del denominador ($\nu_2$)
        \item $F_{1-\alpha, \nu_1, \nu_2} = \frac{1}{F_{\alpha, \nu_2, \nu_1}}$
        \item Usada para comparar varianzas de dos poblaciones normales
    \end{itemize}
\end{tcolorbox}

\begin{multicols}{2}
    \begin{tcolorbox}[title=Tabla de Decisión, equal height group=CH6]
        \begin{center}
        \begin{tabular}{|c|c|c|}
        \hline
        \textbf{Decisión} & \textbf{$H_0$ verdadera} & \textbf{$H_0$ falsa} \\
        \hline
        Rechazar $H_0$ & Error Tipo I ($\alpha$) & Decisión correcta (Potencia $1-\beta$) \\
        \hline
        No rechazar $H_0$ & Decisión correcta ($1-\alpha$) & Error Tipo II ($\beta$) \\
        \hline
        \end{tabular}
        \end{center}
    \end{tcolorbox}
    \begin{tcolorbox}[title=Consideraciones Importantes, equal height group=CH6]
        \textbf{Verificar supuestos}:
        \begin{itemize}
            \item Normalidad (si es necesario)
            \item Tamaño muestral adecuado
            \item Independencia de observaciones
        \end{itemize}
        
        \textbf{Interpretación}:
        \begin{itemize}
            \item "No rechazar $H_0$" $\neq$ "Aceptar $H_0$"
            \item El $p$-valor mide la evidencia contra $H_0$
            \item Considerar el tamaño del efecto, no solo la significación
        \end{itemize}
        
        \textbf{Relación con I.C.}:
        \begin{itemize}
            \item Si $\theta_0$ está fuera del I.C. $(1-\alpha)$: Rechazar $H_0$ al nivel $\alpha$
            \item Si $\theta_0$ está dentro del I.C. $(1-\alpha)$: No rechazar $H_0$ al nivel $\alpha$
        \end{itemize}
    \end{tcolorbox}
\end{multicols}

\end{document}